% Figure: Listing 1, the composable model API in use.
%
% Source for fig_code_compose.pdf / .png. Build it with the sibling driver:
%
%     uv run --group manuscript python manuscript/examples/figures/fig_code_compose.py
%
% The caption is set INSIDE the float and there is deliberately no caption
% paragraph beneath the image in method.md, matching the treatment of
% Algorithm 1 and the venue's practice for code and algorithm boxes.
%
% This code is executed, not illustrative. It was run against the built
% extension on 2026-08-15 and returns r^2 = 0.999198 on a seeded Gaussian, with
% parameters keyed "peak1.sigma", "peak1.center", "peak1.amplitude". If the
% public API changes, this listing is wrong until it is re-run and regenerated.
\documentclass[border=6pt,varwidth=15.2cm]{standalone}

\usepackage[T1]{fontenc}
\usepackage{lmodern}
\usepackage{xcolor}
\usepackage{listings}
\usepackage{caption}

\definecolor{kw}{RGB}{0,90,160}
\definecolor{str}{RGB}{160,60,20}
\definecolor{cmt}{RGB}{110,110,110}
\definecolor{rule}{RGB}{40,40,40}

\lstset{
  language=Python,
  basicstyle=\ttfamily\small,
  keywordstyle=\color{kw}\bfseries,
  stringstyle=\color{str},
  commentstyle=\color{cmt}\itshape,
  showstringspaces=false,
  columns=fullflexible,
  keepspaces=true,
  breaklines=false,
  aboveskip=0.6em,
  belowskip=0pt,
  % Lets $...$ render as math inside the listing, which is how R^2 in the
  % trailing comment becomes a real superscript rather than a caret. Safe
  % because this snippet contains no literal '$'; adding one would need
  % escapeinside instead.
  mathescape=true,
}

\captionsetup{labelfont=bf,skip=0.4em,justification=justified,
              singlelinecheck=false,font=small}
\renewcommand{\lstlistingname}{Listing}

\begin{document}
\begin{minipage}{15.2cm}
\hrule height 0.9pt
\vspace{0.5em}
\captionof{lstlisting}{One model, composed and fitted. The caller states the
model as a graph of named components and names a solver, or asks for
\texttt{auto} to have the solver family derived from the graph; the parameter
layout follows from the graph, and each kernel carries its own Jacobian,
closed-form for all but eleven. The same construction scales to several
components, tied parameters, stacked datasets and N-dimensional fits without a
change of interface.}
\vspace{0.2em}
\hrule height 0.4pt
\begin{lstlisting}
from spectrafit_core import MeasurementData, compose, fit, gaussian

# x, y are the measured abscissa and ordinate.
graph = compose([gaussian("peak1", amplitude=1.0, center=0.0, sigma=1.0)]).build()
result = fit(graph, MeasurementData(x=list(x), y=list(y)))

result.parameters   # keyed "peak1.<param>"; each carries .value and .stderr
result.r_squared    # goodness of fit, as $R^2$ against the data
\end{lstlisting}
\vspace{0.3em}
\hrule height 0.9pt
\end{minipage}
\end{document}
